\SGAProof{Proof}\footnote{Editor's note: the proof has been expanded.}
Write \(\pi\colon G\to S\) for the structure morphism and
\(\epsilon\colon S\to G\) for the unit section.  Since \(G\) is flat
over \(S\) at \(\epsilon(s_0)\), and since \(G_{\bar s_0}\) is reduced
(hence smooth over \(\bar s_0\), by VI\textsubscript{A}, \SGARefLink{prop:sga3-VIA-1.3.1}{1.3.1}),
\(G\) is smooth over \(S\) at \(\epsilon(s_0)\)
(EGA IV\textsubscript{4}, 17.5.1).  Thus there is an open
neighborhood \(V\) of \(\epsilon(s_0)\) on which \(\pi\) is smooth.
Then \(S'=\epsilon^{-1}(V)\) is open in \(S\), and \(G_{S'}\) is
smooth over \(S'\) at the points of \(\epsilon(S')\).  Since \(G\)
has connected fibers, VI\textsubscript{B}, \SGARefLink{thm:sga3-VIB-3.10}{3.10} shows that \(G_{S'}\)
is smooth over \(S'\).  Replacing \(S\) by \(S'\), we may therefore
suppose that \(G\) is smooth over \(S\).

By \SGARefLink{sga3:xi:theorem:4:1}{Exposé XI, Theorem 4.1}, the functor of subgroups of multiplicative
type of \(G\) is represented by an \(S\)-scheme \(\mathscr M\), smooth
and separated over \(S\).  Let \(r_0\) be the rank of
\(G_{\bar s_0}\), and consider the open subscheme
\(\mathscr M_{r_0}\subset\mathscr M\) representing the functor of
subtori of rank \(r_0\) of \(G\) (IX, \SGARefLink{sga3:ix:definition:1:4}{1.4}).  Smoothness implies that
\(\mathscr M_{r_0}\) has a rational point over a finite separable
extension of \(\kappa(s_0)\) (EGA IV\textsubscript{4}, 17.15.10(iii)).
Consequently there is an étale morphism \(S'\to S\), together with a
point \(s'_0\) above \(s_0\), such that \(G_{s'_0}\) has a torus of
rank \(r_0\).  Replacing \(S\) by \(S'\), we may therefore suppose
that \(G_{s_0}\) has a torus of rank \(r_0\).

By ``Hensel's lemma'' (XI, \SGARefLink{sga3:xi:proposition:1:10}{1.10}), the smoothness of
\(\mathscr M_{r_0}\) allows this torus to be lifted to an \(S'\)-torus
\(T\) of \(G\), where \(S'\to S\) is étale and is equipped with a
point \(s'_0\) above \(s_0\).  By \SGARefLink{sga3:x:corollary:4-5}{Exposé X, 4.5}
(see also \SGARefLink{sga3:xix:proposition:6-1}{6.1}\footnote{Editor's note: §\SGARefLink{sga3:xix:proposition:6-1}{6.1} is independent of the rest
of this Exposé.}), there is an étale morphism \(f\colon S''\to S'\)
that splits \(T\) and for which \(f^{-1}(s'_0)\ne\varnothing\).
Because an étale morphism is open, and the assertions in \textup{\SGARefLink{sga3:xix:theorem:2:5:item:i}{(i)}}
are local for the étale topology, we may thus suppose that \(G\) has
a split torus \(T\),\footnote{Editor's note: throughout this volume,
``trivial torus'' has been replaced by ``split torus''.} maximal at
\(s_0\).  Write \(T=D_S(M)\), and let
\[
  \mathfrak g=\prod_{m\in M}\mathfrak g^{m}
\]
be the decomposition of
\(\mathfrak g=\operatorname{Lie}(G/S)\) under
\(\operatorname{Ad}(T)\) (\SGARefLink{sga3:i:proposition:4-7-3}{Exposé I, 4.7.3}).

Put \(\mathfrak t=\operatorname{Lie}(T)\), and, for every \(m\in M\),
write
\[
  \mathfrak g^{m}(s_0)
  =
  \mathfrak g^{m}\otimes_{\mathcal O_S}\kappa(s_0).
\]
\footnote{Editor's note: the preceding sentence has been added.}
Let \(R\) be the set of \(m\in M\), \(m\ne0\), such that
\(\mathfrak g^{m}(s_0)\ne0\).\footnote{Editor's note: since
\(\mathfrak g\) is a finite locally free \(\mathcal O_S\)-module,
\(R\) is finite.}  Since \(G_{\bar s_0}\) is reductive,
\(\mathfrak g^{0}(s_0)=\mathfrak t(s_0)\), and hence
\[
  \mathfrak g(s_0)
  =
  \mathfrak t(s_0)
  \oplus
  \prod_{\alpha\in R}\mathfrak g^{\alpha}(s_0).
\]
Since the modules in question are locally free, after shrinking \(S\)
we may suppose that the \(\mathfrak g^{\alpha}\) are free and that
\[
  \mathfrak g
  =
  \mathfrak t
  \oplus
  \prod_{\alpha\in R}\mathfrak g^{\alpha}.
\]

Recall from \SGARefLink{sga3:xii:1.12}{Exposé XII, 1.12} that a group of multiplicative type has a
unique maximal torus (this also follows almost immediately by descent,
the diagonalizable case being evident).  Let \(T_\alpha\) be the
maximal torus of \(\ker(\alpha)\), and put
\[
  Z_\alpha=\underline{\operatorname{Centr}}_G(T_\alpha).
\]
\footnote{Editor's note: the following argument has been expanded.}
By \SGARefLink{sga3:xix:paragraph:2-2}{2.2}, \(Z_\alpha\) is affine and smooth over \(S\), with connected
fibers; and by \SGARefLink{sga3:xix:proposition:1-12}{1.12}, its fiber
\((Z_\alpha)_{\bar s_0}\) is reductive, of semisimple rank \(1\), with
maximal torus \(T_{\bar s_0}\).  By \SGARefLink{sga3:xix:lemma:2-3}{2.3} there is therefore an open
subset \(U_\alpha\subset S\), containing \(s_0\), on which
\(Z_\alpha\) has reductive fibers.

Set
\[
  U=\bigcap_{\alpha\in R}U_\alpha.
\]
By \SGARefLink{sga3:xix:proposition:1-12}{1.12}, for every \(s\in U\), \(G_{\bar s}\) is reductive, with
maximal torus \(T_{\bar s}\), and the set of roots of \(G_{\bar s}\)
relative to \(T_{\bar s}\) is identified with \(R\).  Consequently
\[
  \operatorname{rgred}(G_{\bar s})
  =
  \dim(T)
  =
  \operatorname{rg}(M),
  \qquad
  \operatorname{rgss}(G_{\bar s})
  =
  \operatorname{rg}(R)
  \quad\text{(cf. 1.10).}
\]
This proves \textup{\SGARefLink{sga3:xix:theorem:2:5:item:i}{(i)}} and the first assertion of \textup{\SGARefLink{sga3:xix:theorem:2:5:item:ii}{(ii)}}.
It remains to prove that \(W_{G_U}(T_U)\) is finite over \(U\), that
is, that it has the same number of points in every geometric fiber
(SGA 1, I, 10.9; EGA IV\textsubscript{3}, 15.5.1; and
EGA IV\textsubscript{4}, 8.12.4).

For this it is enough to observe that, for \(s\in U\), the geometric
fiber of this group is determined by the configuration \(R\subset M\):
it is the constant group associated with the abstract Weyl group of
this root system.  In particular it is independent of \(s\); see the
\emph{Bible}, §11.3, Theorem 2 (and also Exposé XXII, no. 3).
\hfill\(\square\)

